% Source: https://artofproblemsolving.com/community/c3h1508944_logic_and__parity
% Stable identity: aops-topic-1508944-q01
\textbf{If room 8 is empty, the puzzle cannot be solved.} With room 8 empty, both of these situations satisfy all the rules: the princess in room 1 and all other rooms empty (sign 1 is true, and no tiger signs need to be false), and the princess in room 7 and all other rooms empty (sign 7 is true). So the king must have said that room 8 is \emph{not} empty.

\textbf{Room 8 is not empty.} Room 8 cannot hold the princess: its sign would be true, but it says that room 8 holds a tiger. So room 8 holds a tiger and sign 8 is false. Its first half (``this room contains a tiger'') is true, so the second half is false: room 9 is not empty.

\textbf{Room 9.} If the princess were in room 9, sign 9 would be true, but it says room 9 holds a tiger. So room 9 holds a tiger, sign 9 is false, and since its first half is true, ``sign 6 is wrong'' is false: \emph{sign 6 is true}.

\textbf{The chain.} Sign 6 true means sign 3 is false, so sign 5 is wrong and sign 7 is right.
\begin{itemize}
\item Sign 7 right: the princess is not in room 1.
\item Sign 5 wrong: sign 2 and sign 4 are both wrong. Sign 4 wrong means sign 1 is right, so the princess is in an odd room: $3$, $5$ or $7$ (not $1$, not $9$). Sign 2 wrong means room 2 is not empty, so it holds a tiger.
\end{itemize}
The princess's sign must be true. Sign 3 is false and sign 5 is wrong, so she is not in room 3 or room 5. She is in room $\boxed{7}$, whose sign (``not in room 1'') is indeed true.

This is consistent: for example, rooms 2, 8, 9 hold tigers (signs 2, 8, 9 false) and rooms 1, 3, 4, 5, 6 are empty.
